%
% Introduction
%

% !TEX root = ../../main.tex

\chapter{Einleitung}

  % TODO nomol..

  Viele Optimierungsprobleme,
  die in der Informatik und Robotik auftreten sind mit mathematischen Methoden nicht innert nützlicher Frist lösbar.
  Oft spielt dabei die Anzahl der Freiheitsgrade eine entscheidende Rolle.
  Die Anzahl der Freiheitsgrade und somit der Menge möglicher Lösungen übersteigt die Kapazität aktueller Computer.
  \\
  \\

  % TODO alternativ zu was

  Ein alternatives Verfahren findet sich in der Natur. In der Natur werden ständig Optimierungsprobleme gelöst.
  Das Prinzip dabei ist, dass die Individuen überleben und Nachwuchs zeugen,
  die am besten an die aktuelle Umgebung angepasst sind.
  Die nächste Generation ist besser an die Bedingungen ihrer Vorgänger angepasst. % ZITAT

  % nicht nur die guten, sondern auch schlechtere

  % kombination aus reproduktion --  kopie --> variation, selektion, vererbung

  % selektions druck, gut reproduzieren sich mit höhere wahrscheinlichkeit
  % führt zu veränderung.

  % 1 Mehr formulieren, dass formale lösungs methoden -> planen
  % 2 natur: trial and error
  %
  % lösungsraum zu kompliziert für planung
  %

  \\
  \\
  Mit dieser Methode können Optimierungsprobleme auf einem Computer gelöst werden,
  da nur ein Teil der ganzen Lösungsmenge berechnet wird.
  Auf Computersystemen kann dieser Prozess virtuell abgebildet und im Schnelllauf durchgeführt werden.

  \section{Ausgangslage}

    \todo[inline]{Nennt bestehende Arbeiten/Literatur zum Thema -> Literaturrecherche}
    \todo[inline]{Stand der Technik: Bisherige Lösungen des Problems und deren Grenzen}

    Evolutionäre Algorithmen sind bereits in vielen Bereichen im Einsatz.
    Besonders bei Problemstellungen wo eine approximative Lösung ausreichend ist,
    wird dieses Verfahren eingesetzt.
    \\
    Die NASA verwendet evolutionäre Algorithmen um Antennen zu entwerfen,
    die im Weltraum zur Funkkommunikation verwendet werden~\cite{Hornby2006}.
    Dabei wurde die Antenne auf gute Sendeleistung bei geringer Energieaufnahme optimiert.
    \\
    Als theoretische Grundlagen wurden die Bücher ``Bio inspired artificial intelligence''~\cite{book:bioInspired} und
    ``Evolutionäre Algorithmen''~\cite{book:evAlgo} verwendet.

    % Altes Problem der Robotik.

  \section{Aufgabenstellung}

    % TODO aufgabenstellung zitieren.

    % nicht fkt der zeit, sondern zeit

    Als Mass für die Fitness wird eine Funktion der Zeit verwendet, welche das Tier braucht,
    um einen gegebenen Parcours zurückzulegen.

    % satz weg

    Dabei müssen Individuen berücksichtigt werden, die den Parcours nicht vollständig ablaufen können.
    Die Fitnessfunktion berücksichtigt in diesem Fall die zurückgelegte Strecke.

    % klarer formulieren, wie es ist.
    \\
    Als Randbedingungen vorgegeben sind:
    \begin{enumerate}
      \item Forderung einer physikalisch sinnvollen Bewegung
      \item Obere Grenze für die total aufgewandte Energie
      \item Beschränkung der abgegebenen Leistung
    \end{enumerate}
    Damit die Forderung nach einer physikalischen Bewegung erfüllt werden kann,
    wird die Bewegung mit Hilfe einer \gls{PhysicsEngine} simuliert.
    Diese Engine wird von den Studierenden ausgewählt und kann als Blackbox verwendet werden.

  \section{Zielsetzung}

    % TODO holprig
    % reihenfolge !

    % !!! wissenschaftliche fragenstellung formulieren !!!
    % anzahl (1,2,55) klare fragen stellen: gesucht ist ein tier, welche den parcours in endlicher zeit bewältigt.
    % 2. welche anzahl beine ist optimal
    %...
    % möglich, dass auf spätere hypothesen verwiesen wird

    Das Ziel dieser Arbeit ist sechsbeinigen künstlichen Tieren mit Hilfe eines evolutionären Algorithmus,
    in einer zweidimensionalen Umgebung das Gehen zu beizubringen.
    Um dieses Ziel zu erreichen, soll eine Applikation erstellt werden,
    welche den evolutionären Algorithmus implementiert.
    Die Applikation muss ebenfalls die Möglichkeit bieten, eine Simulation durchzuführen.
    Es soll nicht nur die Steuerung der Beine evolviert werden (Gehen lernen),
    sondern auch die Form des Körpers und die Abmessung der Beine.
    Zur Implementierung soll eine bestehende \gls{PhysicsEngine} verwendet werden.
    Der Hauptfokus liegt auf dem Evolvieren der artifiziellen Tiere und nicht auf dem Ausbau der \gls{PhysicsEngine}.
    Zusätzliche Hypothesen~(\vref{subsub:hypoKp} und~\vref{subsub:hypoMut}) werden im Kapitel Methode~(\vref{sec:method}) erwähnt,
    da die erforderlichen Begrifflichkeiten dann eingeführt sind.

  \section{Anforderungen}

    \subsection{Fitnessfunktion}

      Die Fitness ist ein Mass für den Grad der Anpassung eines Individuums an die Umgebung.
      In frühen Generationen können Tiere existieren, die sich nicht fortbewegen können.
      Die Zeit, welche ein Individuum für die Bewältigung eines Parcours (\vref{subsub:parcours}) brauchen würde,
      wäre somit unendlich.
      Um dieses Problem zu umgehen, wird als Mass für die Fitness die zurückgelegte Strecke verwendet.
      Für die Bewältigung des Parcours steht dem Individuum eine fixe Anzahl Zeiteinheiten zur Verfügung.

    \subsection{Parcours\label{subsub:parcours}}

      Die Individuen müssen während der Simulation eine Strecke zurücklegen.
      Diese Strecke und ihre Begebenheiten werden als Parcours bezeichnet.
      \vref{fig:scetchParcours} zeigt eine Skizze des Parcours.
      Mit zunehmenden Generationen soll die Schwierigkeit des Parcours steigen.
      Als Schwierigkeit eines Parcours wird die maximale Steigung bezeichnet.
      Der Parcours soll zufällig generiert werden.

      \begin{figure}[H]
        %
% Concept parcours
%

% !TEX root = ../main.tex

\begin{tikzpicture}[scale=\textwidth/20cm]

  \coordinate (P0) at (0,0);
  \coordinate (P1) at (0,1);
  \coordinate (P2) at (4,1);
  \coordinate (P3) at (7,2);
  \coordinate (P4) at (9,2);
  \coordinate (P5) at (10,1.5);
  \coordinate (P6) at (12,3);
  \coordinate (P7) at (14,3);
  \coordinate (P8) at (16,2.2);
  \coordinate (P9) at (20,2.2);
  \coordinate (P10) at (20,0);

  % Draw parcours
  \draw [fill=black] (P0) -- (P1) -- (P2) -- (P3) -- (P4) -- (P5) -- (P6) -- (P7) -- (P8) -- (P9) -- (P10) -- cycle;

\end{tikzpicture}

        \caption{Skizze Parcours\label{fig:scetchParcours}}
      \end{figure}

    \subsection{Limiten\label{sub:IntroReqLimit}}

      Limiten verhindern, dass der Parcours durch triviale Lösungen wie z.B. unendlich grosse Beine bewältigt werden kann.
      Ebenfalls schränken sie den Lösungsraum ein und helfen gewinnbringende Lösungen zu finden.
      Für folgende Komponenten werden Limiten eingeführt:
      % Warum dieser limiten?
      \begin{itemize}
        \item Winkelgeschwindigkeit auf den Gelenken: Verhindert Beschleunigung ins unendliche.
        \item Beinhöhe: Verhindert die triviale Überwindung von Hindernissen.
        \item Beinbreite: Verhindert unbrauchbare Lösungen
        \item Masse des gesamten Individuums: Vereinfachung der Problemstellung
        \item Fläche des Körpers: Verhindert unendlich grosse Individuen
      \end{itemize}

    \section{Projektmanagement}

      % TODO in anhang verschieben?

      Die verschiedenen Arbeitspakete werden eine Woche im Voraus geplant.
      Dieses Vorgehen wird auch als rollende Planung bezeichnet.
      Ein Projektvorgehensmodell wird nicht benutzt,
      jedoch ist eine Projektplanung mit Hilfe von MS Project 2016 erstellt worden.
      Alle wichtigen Meilensteine finden sich im Projektplan.
      Der Stand der Arbeit wurde an einem wöchentlichen Meeting mit den Betreuern besprochen.
      Das Beschlussprotokoll wurde vor dem nächsten Meeting verschickt.
